\chapter{The Largest Prize: Why AI for Science Dominates the Value Case}
\chaptermark{The Largest Prize}
\label{ch:prize}

This book has so far treated AI as an industrial commodity and asked who will supply it. This chapter asks what it is \emph{for}, and argues a claim that most of the public debate has inverted: the largest value in artificial intelligence is not in writing software, drafting documents, or automating enterprise workflow. It is in accelerating science. Everything else redistributes an existing surplus; science creates new ones. If that ranking is right, then the two chapters that follow---how much compute a nation's research actually needs, and what institution could supply it---are not a digression from this book's argument but its destination, and the competition described in Parts~I--III is a competition over who gets to compound their scientific output fastest.

The chapter proceeds in the order that honesty requires: first the demonstrated results, then an unsparing ledger of what has \emph{not} worked, then the economics that make the prize large even after the failures are subtracted, then the strategic consequence.

\section{What has actually been demonstrated}

Four domains have produced results that survive scrutiny, and their common feature is instructive: in each, AI replaced a computation that was previously bounded by physics-based simulation cost, and the replacement was validated against an established operational baseline rather than against a benchmark of its own devising.

\textbf{Weather is the strongest case, and the only one with operational adoption by three national agencies.} ECMWF's AIFS became operational in February 2025 at 28\,km grid spacing with gains up to 20\% over the physics-based model on measures including tropical-cyclone tracks, at roughly \emph{one-thousandth} of the energy per forecast; its 51-member ensemble followed in July 2025~\cite{ecmwf-aifs}. NOAA deployed AI-driven global models operationally in December 2025: a 16-day forecast in about 40 minutes using \emph{0.3\%} of the compute of its conventional system, with the ensemble variant extending useful skill by 18--24 hours~\cite{noaa-ai}. Microsoft's Aurora---1.3B parameters, pretrained on 32 GPUs in about two and a half weeks---beats the operational high-resolution model on 92\% of targets and matched or beat all seven operational centres on every tested cyclone track, at roughly $10^5$ times the speed of the atmospheric-chemistry system it replaces; adapting it to a new task takes four to eight weeks against the several years a new dynamical model requires~\cite{aurora}. NeuralGCM runs on a single accelerator what previously required 13,000 CPUs~\cite{neuralgcm}. China's meteorological administration has its own operational family, and a 2026 typhoon-track system reports up to 32\% error reduction against leading global ensembles~\cite{cma-ai,fuxi-cnops}.

\textbf{Structural biology is the most diffused case.} AlphaFold has predicted over 200 million structures, is used by more than three million researchers in 190 countries---including over a million in low- and middle-income countries---and won the 2024 Nobel Prize in Chemistry~\cite{alphafold-impact,nobel-2024}. The measurable second-order effect matters more than the structures: researchers using it increased their submissions of \emph{novel experimental} structures by over 40\% and were twice as likely to be cited in clinical work~\cite{alphafold-impact}. The database now includes millions of predicted complexes, representing on the order of 17 million GPU-hours of computation given away~\cite{alphafold-db}. Protein \emph{design} has followed: a generative model produced a fluorescent protein 58\% identical to its nearest natural relative---an object evolution had not found---and diffusion-based enzyme design has produced metallohydrolases with catalytic efficiency orders of magnitude above previous designed enzymes~\cite{esm3,rfdiffusion3}.

\textbf{Mathematics and algorithm discovery produced the first credible autonomous results.} A frontier system achieved gold-medal standard at the 2025 International Mathematical Olympiad, solving five of six problems end-to-end in natural language within the human time limit~\cite{imo-gold}. More consequentially for this book, an evolutionary coding agent found a 48-multiplication algorithm for $4\times4$ complex matrix multiplication---the first improvement on that case since Strassen in 1969---improved the best known solution on 20\% of fifty-plus open problems it was set, and returned measurable production value: 0.7\% of a hyperscaler's worldwide compute recovered through better scheduling, a 23\% kernel speedup yielding a 1\% reduction in the training time of the company's own frontier model~\cite{alphaevolve}. That last result is the flywheel of Section~\ref{sec:flywheel} in miniature: AI improving the infrastructure that trains AI.

\textbf{Drug discovery has its first controlled clinical datapoint.} An AI-discovered TNIK inhibitor for idiopathic pulmonary fibrosis reported Phase~IIa results in \emph{Nature Medicine} in June 2025---71 patients, 12 weeks, a $+98.4$\,mL mean forced-vital-capacity change against $-20.3$\,mL for placebo---and entered a 320-patient Phase~III in July 2026; the developer reports reaching preclinical candidate nomination in 12--18 months with 60--200 molecules synthesized per programme, against a conventional 2.5--4 years~\cite{rentosertib,rentosertib-p3}. This is one drug, one surrogate endpoint, one country, and not an approval. It is also more than the field had two years ago.

\section{The failure ledger}

A chapter arguing that AI4S is the largest prize is obliged to lead with the evidence against itself, and in this field the counterevidence is unusually sharp.

\emph{The most-cited economic study of AI in scientific discovery was fabricated.} A 2024 paper reporting large gains in materials discovered, patents, and product innovation at a real firm was repudiated by its institution in May 2025---``no confidence in the provenance, reliability or validity of the data''---and the two prominent economists who had publicly praised it jointly stated that its findings should not be relied upon~\cite{mit-repudiation}. It is not cited here except as a warning, and any AI4S argument that still rests on it should be discarded.

\emph{Materials discovery, the flagship domain, has not survived validation intact.} The autonomous-synthesis result that anchored the field was formally corrected in January 2026---successes cut from 40 to 36, one compound removed because it was in the training data, and the authors conceding the materials were ``new to the prediction platform, not necessarily new to science''~\cite{alab-correction}. Independent crystallographers argued that none of the originally reported compounds was genuinely new~\cite{alab-critics}. The 380,000-stable-crystal result has been criticized by senior materials chemists for producing ``chemical compounds rather than materials'' with no demonstrated functionality~\cite{gnome-critique}. The strongest generative-materials paper validated exactly one compound experimentally, and it missed its target modulus by about 20\%~\cite{mattergen}.

\emph{Even weather has a documented failure mode.} In May 2026 ECMWF discontinued real-time running of all four external AI models, having found them increasingly sensitive to the exact initializing analysis---several degraded when initialized on an upgraded analysis system---and noting that two of them do not forecast precipitation at all~\cite{ecmwf-farewell}. The lesson is not that AI weather models failed; the centre's own AIFS remains operational. It is that these models are couplings to an underlying data-assimilation system, not free-standing oracles.

\emph{Autonomous discovery is far weaker than its marketing.} When a frontier agent was run against 700 open problems in one mathematical corpus, human adjudication of 200 candidate solutions found 68.5\% fundamentally flawed and only 6.5\% meaningfully correct---of which two were genuine autonomous resolutions~\cite{aletheia}. A leading mathematician's assessment of a widely publicized single success was that only 1--2\% of the open problems in that corpus are tractable this way, that difficulty within it varies by orders of magnitude, and that unreported failures skew the perceived rate~\cite{tao}. The first conference to accept AI as primary author found submissions ``technically correct'' but ``neither interesting nor important,'' with reviewers noting that technical fluency can mask poor scientific judgement~\cite{agents4science}. Working scientists interviewed about frontier models in early 2026 were blunt: ``We have not found, yet, that LLMs are fundamentally changing the way that science is done''~\cite{openai-science}.

\emph{And the scientific commons is absorbing real damage.} One journal measuring its own pipeline found submissions up 42\% since late 2022, over 30\% of peer reviews showing detectable AI use, and papers with heavy AI content receiving revise-and-resubmit at 3.2\% against 11.9\% for low-AI papers---while the editorial board nearly doubled to cope~\cite{ai-slop}.

The honest summary: AI4S has delivered decisively where the task is \emph{a well-posed prediction problem with an expensive physics-based incumbent and an established validation baseline}, and has so far mostly failed where the task is \emph{judgement}---choosing which question is worth asking, deciding what counts as new, or assessing whether a result matters. That boundary is the most useful thing known about the field, and it has a direct architectural implication taken up in Section~\ref{sec:prize-arch}.

\section{Why the prize is nevertheless the largest one}

\subsection{The denominator}

Start with the sizes of the things being compared. Global research and development spending ran approximately \$3.1 trillion (PPP) in 2022, with the United States at \$923B and China at \$812B~\cite{nsf-indicators}. The entire worldwide software market---the addressable universe of ``AI writes code'' and ``AI streamlines the enterprise''---is forecast at \$1.47 trillion for 2026, roughly half of global R\&D~\cite{gartner}. And the combined revenue of the frontier laboratories whose valuations Chapter~\ref{ch:trade} examined is under 2\% of annual global R\&D spending~\cite{lab-revenue}. Even the most-cited bullish estimate of generative AI's economic potential put \emph{R\&D as the single largest affected function} at about a quarter of the total, ahead of software engineering---while explicitly excluding value from entirely new product categories, which is precisely where scientific advance shows up~\cite{mckinsey-genai}.

\subsection{The return}

Research spending is not an ordinary input. The best available estimate of the social return to R\&D is \$13.3 of social benefit per dollar spent at a 5\% discount rate---and even under a deliberately punitive 20-year-lag assumption, \$4.9 per dollar; the authors note it is difficult to construct an average return below \$4~\cite{jones-summers}. No enterprise-software deployment has a multiplier in that range, because enterprise software largely reallocates surplus between firms while research generates surplus that is mostly not captured by whoever paid for it. That non-appropriability is exactly why public institutions fund research, and exactly why the AI4S case is a \emph{public-infrastructure} case rather than a product case---a point Chapter~\ref{ch:consortium} builds an institution on.

Three accounting disciplines keep this subsection honest. \emph{First}, the \$4--13 figure is an \emph{intertemporal present value} of the social surplus stream generated by a dollar of R\&D, not an annual multiplier; it must never be multiplied directly against a single year's R\&D flow to produce an annual benefit. The correct chain is: an annual \emph{productivity uplift} to the research enterprise (the funnel's output, below) $\rightarrow$ an increment to the effective R\&D flow $\rightarrow$ a \emph{present value} of social surplus obtained by applying the Jones--Summers ratio to that increment. The three quantities---annual uplift, incremental output flow, and present value of surplus---are kept separate wherever this book uses them. \emph{Second}, the ratio is an \emph{average} return; AI will first reach the most computationally tractable tasks, whose marginal return may sit above or below that average, and nothing in the current evidence pins down which. \emph{Third}, the \$3.1T denominator is PPP-adjusted while software and laboratory revenues are nominal; the comparison is strategically suggestive, not an accounting identity. None of these disciplines shrinks the qualitative conclusion---the base is enormous and the multiplier is large---but they change what may legitimately be computed from it.

\subsection{The bottleneck}

Now the argument that makes AI4S strategic rather than merely valuable. Research productivity is falling. The canonical measurement finds aggregate US research productivity down by a factor of 41 since the 1930s---more than 5\% per year---sustained only by a 23-fold increase in research effort; keeping Moore's law required more than eighteen times as many researchers per doubling as in the early 1970s; cancer clinical-trial research productivity fell by a factor of 4.8 between 1975 and 2006~\cite{ideas-harder}. Ideas are getting harder to find, and the traditional remedy---hire more researchers---is exhausting itself demographically in every advanced economy.

An important 2026 counter-interpretation deserves its place here, because it redirects rather than refutes the argument. Fort and coauthors find that patent production per unit of R\&D input may \emph{not} be declining at all; what has weakened is the conversion of ideas into firm growth---the diffusion step, not the discovery step~\cite{fort2026}. The distinction matters strategically. If ideas are genuinely harder to produce, AI's leverage is on discovery itself, and the funnel below is the right model. If ideas are produced but diffuse poorly, the binding constraint is institutional---standards, competition, technology transfer, absorptive capacity---and the highest-return intervention is the deployment and certification infrastructure of Chapters~\ref{ch:trust} and~\ref{ch:consortium} rather than more model capability. The truth almost certainly varies by sector, which is one more reason the consortium proposal pairs compute provision with diffusion machinery instead of treating capability as self-executing. Either way the case for the investment survives; what changes is where the marginal dollar goes.

An input that raises research productivity therefore acts on the binding constraint of long-run growth, which is why the economics of AI turns on this question and not on call-centre automation. And here is the pivot of this chapter. The most influential \emph{skeptical} estimate of AI's macroeconomic effect---a total-factor-productivity gain of about 0.66\% over ten years, roughly one-tenth of bullish estimates---reaches that number by counting only task-level automation of existing work, and \emph{explicitly excludes} transformative science-driven gains on the stated grounds that advances of the protein-folding type ``do not seem likely within the 10-year time frame''~\cite{acemoglu}. That exclusion is the entire disagreement. Structure-prediction, operational AI weather, and designed enzymes all arrived inside the window. One does not have to accept the most aggressive growth models---one recent analysis finds that automating as little as 13\% of research tasks economy-wide is sufficient to enter an explosive-growth regime~\cite{davidson}---to notice that the credible range of estimates for AI's contribution to annual growth spans from 0.1\% to 30\%, and that the single variable separating those poles is whether AI automates \emph{innovation itself}~\cite{quadrillion}.

\subsection{The flywheel}
\label{sec:flywheel}

The final reason AI4S outranks the alternatives is that it is the only application category that feeds back into the means of its own production. Better materials improve batteries, magnets, and semiconductors; better semiconductor design improves accelerators; better weather and grid modelling improves the energy system that Chapter~\ref{ch:energy} identified as the binding physical constraint; better algorithms improve training efficiency directly, as the 1\%-of-training-time result cited above demonstrates in the smallest possible way. Each of these is an input to the war this book describes. Enterprise-workflow automation has no such loop: a faster expense-report pipeline does not lower the cost of the next accelerator generation.

\section{The value-conversion funnel}
\label{sec:funnel}

The claim that AI4S is the largest prize needs a conversion model, not just a denominator, because value does not pass unchanged from capability to output. The chain runs
\[
\text{AI capability} \rightarrow \text{task acceleration} \rightarrow \text{validated discovery} \rightarrow \text{deployment} \rightarrow \text{productivity},
\]
and each arrow has a conversion efficiency well below one. Stating them explicitly---even as ranges---is the difference between an argument and an enthusiasm.

\begin{figure}[htbp]
\centering
\begin{tikzpicture}[
  scale=0.86, transform shape,
  every node/.style={font=\scriptsize},
  st/.style={draw=inkMut, line width=0.6pt, rounded corners=2pt, align=center,
             inner sep=4pt, minimum height=13mm},
  ar/.style={-{Stealth[length=5pt]}, line width=1pt, color=inkSec},
  eff/.style={font=\scriptsize\color{sTwo}, align=center},
]
\node[st, fill=sOne!12, draw=sOne, text width=26mm] (a) at (0,0) {AI capability};
\node[st, fill=sOne!10, draw=sOne, text width=24mm, right=11mm of a] (b) {task\\acceleration};
\node[st, fill=sThree!10, draw=sThree, text width=23mm, right=11mm of b] (c) {validated\\discovery};
\node[st, fill=sTwo!10, draw=sTwo, text width=21mm, right=11mm of c] (d) {deployment};
\node[st, fill=sFour!12, draw=sFour, text width=20mm, right=11mm of d] (e) {productivity};

\draw[ar] (a) -- node[above, eff] {15--35\%\\addressable} node[below, font=\scriptsize\color{inkSec}] {} (b);
\draw[ar] (b) -- node[above, eff] {1.5--5$\times$\\on that part} (c);
\draw[ar] (c) -- node[above, eff] {20--60\%\\survives} (d);
\draw[ar] (d) -- node[above, eff] {30--70\%\\converts} (e);

\node[font=\scriptsize\color{inkSec}, align=center, text width=30mm, below=2mm of a] {measured by\\benchmarks};
\node[font=\scriptsize\color{inkSec}, align=center, text width=26mm, below=2mm of b] {weather: $10^3$--$10^5\times$\\structure: 200M};
\node[font=\scriptsize\color{inkSec}, align=center, text width=26mm, below=2mm of c] {materials: attrition\\documented};
\node[font=\scriptsize\color{inkSec}, align=center, text width=24mm, below=2mm of d] {weather $<$1\,yr;\\pharma $\sim$10\,yr};
\node[font=\scriptsize\color{inkSec}, align=center, text width=24mm, below=2mm of e] {\$3.1T base,\\\$4--13 multiplier};

\node[font=\scriptsize\color{inkSec}, align=left, text width=168mm, anchor=north west] at (-2.6,-3.5)
  {The corner-to-corner product of these ranges spans $U=a(s{-}1)vd \approx 0.45\%$ to $59\%$; the task-weighted
   Monte Carlo of \S\ref{sec:funnelmath}, which is the operative estimate, concentrates this into a central
   $\approx$11\% with a 10--90\% range of 8--18\% [M]. Even the pessimistic corner is a socially consequential
   annual research-output flow ($\sim$\$28B/yr against the \$3.1T base), whose present value exceeds the cost of
   national AI4S infrastructure by a large multiple. What the claim does not survive is the assumption that
   acceleration equals discovery, which is why the funnel is drawn rather than asserted.};
\end{tikzpicture}
\caption[The AI4S value-conversion funnel]{From capability to output, with the losses shown. Each arrow carries a
conversion efficiency well below one, and the domains where AI4S value is already visible in the statistics
(weather) are precisely those where all four conversions are favourable---high addressability, a computational
bottleneck, an established validation baseline, and short deployment lag. Where any one fails, as in materials
synthesis, headline discovery counts do not become value. Ranges are author estimates [M]; the purpose is to make
the argument falsifiable term by term rather than to claim precision.}
\label{fig:funnel}
\end{figure}


Four disciplines govern the funnel. \emph{First, addressability}: not all of the \$3.1T R\&D denominator is reachable by current or near-term AI. The reachable share is dominated by the task classes where the evidence above is strong---prediction with an established validation baseline, and search or optimization over well-defined spaces---and is much thinner for autonomous experimental planning and thinner still for scientific judgement and theory formation. A defensible current addressable share is on the order of 15--35\% of research \emph{effort}, weighted toward computational and screening-heavy fields [M]. \emph{Second, acceleration is not discovery}: a hundred-thousand-fold speedup in a forward model, as in weather, converts into scientific value only where the bottleneck was that computation; where the bottleneck is experimental throughput, funding, or judgement, the speedup converts at a small fraction. \emph{Third, validation attrition is severe and measurable}: the materials record of the previous section is the cautionary case, where headline discovery counts survived external scrutiny at rates well below one. \emph{Fourth, deployment and regulatory lag} truncate everything: a validated drug candidate is a decade from revenue, a validated material perhaps half that, a validated forecast improvement is operational within a year---which is why weather is the domain where AI4S value is already visible in the statistics and pharmaceuticals is the domain where it is not yet.

\subsection{The arithmetic, done properly}
\label{sec:funnelmath}

It is often said that multiplying these ranges yields a 3--30\% uplift. It does not, and the error is worth walking through, because the corrected arithmetic is what the rest of this section rests on. Define the output variable precisely: the \emph{incremental} uplift to realized research output from AI, for a single homogeneous task pool, is
\begin{equation}
U = a\,(s-1)\,v\,d,
\label{eq:funnel-simple}
\end{equation}
whose four factors are the four arrows of the funnel: $a$ is \emph{addressability}, the share of research effort that current AI can reach at all; $s$ is the \emph{speed-up} on that reachable share, so that $(s-1)$ counts only the acceleration \emph{beyond} the unassisted baseline (a model that changes nothing has $s=1$ and contributes zero); $v$ is \emph{validation survival}, the fraction of accelerated output that withstands independent scrutiny; and $d$ is \emph{deployment conversion}, the fraction of validated results that actually reach use. All four are dimensionless fractions except $s$, which is a multiplier. The printed endpoint ranges then give $U_{\min}=0.15\times0.5\times0.20\times0.30=0.45\%$ and $U_{\max}=0.35\times4\times0.60\times0.70=58.8\%$; the gross converted output $U^{\mathrm{gross}}=asvd$ spans 1.35--73.5\%. The corner-to-corner interval is therefore roughly \emph{half a percent to sixty percent}, not 3--30\%, and the wider band is the honest one: it is what independence of the terms actually implies.

The corners, however, are not equally credible, because the parameters are correlated in a specific, evidence-backed way: domains with high acceleration and high addressability (weather, structure prediction) also tend to have strong validation infrastructure and short deployment lags, while domains with weak validation (materials synthesis) also carry long deployment lags. The all-pessimistic and all-optimistic corners are thus doubly unlikely. The right formulation---and the one adopted here---is task-weighted:
\begin{equation}
U_{\mathrm{AI4S}} \;=\; \sum_j w_j\, a_j\,(s_j - 1)\, v_j\, d_j\, \ell_j,
\label{eq:funnel-weighted}
\end{equation}
where $j$ indexes task classes, $w_j$ is each class's share of global research effort, and $\ell_j$ is a lag-and-friction discount absorbing calendar delay, organizational adjustment, crowd-out, and regulatory latency---the term the single-pool version silently set to one. Table~\ref{tab:funnel-classes} states the class-level judgments [M] that the evidence of this chapter supports; this ``jagged scientific production function''---strong augmentation of search and prediction, weak automation of judgment---is the same structure the emerging economics literature on AI in science models formally~\cite{agrawal2026}.

\begin{table}[htbp]
  \centering
  \scriptsize
  \caption{The funnel by task class [M], now fully reproducible: every class prints its complete central tuple $(w_j, a_j, s_j, v_j, d_j, \ell_j)$ and the contribution $\Delta U_j = w_j a_j (s_j{-}1) v_j d_j \ell_j$ computed from it exactly. Pessimistic and optimistic values for every parameter, and the formulas, are published in the machine-readable register (\texttt{funnel-model.csv}). Weights sum to one.}
  \label{tab:funnel-classes}
  \setlength{\tabcolsep}{4.5pt}
  \begin{tabular}{p{33mm}ccccccc}
    \toprule
    Task class & $w_j$ & $a_j$ & $s_j$ & $v_j$ & $d_j$ & $\ell_j$ & $\Delta U_j$ \\
    \midrule
    Prediction \& surrogates & 0.20 & 0.50 & 3.0 & 0.70 & 0.70 & 0.70 & 6.9\% \\
    Search \& optimization & 0.20 & 0.40 & 3.5 & 0.40 & 0.60 & 0.60 & 2.9\% \\
    Experimental planning & 0.20 & 0.25 & 2.5 & 0.40 & 0.60 & 0.50 & 0.9\% \\
    Interpretation \& causal inference & 0.25 & 0.15 & 1.5 & 0.30 & 0.50 & 0.35 & 0.1\% \\
    Question selection \& theory & 0.15 & 0.05 & 1.5 & 0.20 & 0.30 & 0.20 & 0.0\% \\
    \midrule
    \textbf{Central total} & & & & & & & \textbf{10.7\%} \\
    \bottomrule
  \end{tabular}
\end{table}

A table of this kind fails silently if it omits the deployment column $d_j$---an omission that sets deployment conversion to one in the printed contributions without ever saying so---which is why the discipline here is that every published number be generated from published tuples. Table~\ref{tab:funnel-classes} complies: each $\Delta U_j$ is the exact product of its row, the central total is their exact sum (10.7\%), and the pessimistic/central/optimistic triple for every parameter lives in a machine-readable CSV shipped with the manuscript source, so any reader can recompute or re-weight the model in one spreadsheet.

The Monte Carlo behind that table is run rather than promised. Sampling each parameter from a triangular distribution over its (pessimistic, central, optimistic) triple, with the three conversion terms $(v_j, d_j, \ell_j)$ coupled within and across classes through a common ``conversion-quality'' factor (mixing weight 0.5)---the correlation structure the evidence supports, since domains with good validation also deploy faster---200{,}000 draws give: \emph{median 11.8\%, 10th--90th percentile 7.7--17.5\%} [M]. One-at-a-time sensitivity ranks the classes: the prediction-and-surrogates class dominates (swinging the total from 4.6\% to 42.9\% across its own corner tuples), search-and-optimization second (8.0--27.7\%), and the two judgment-heavy classes barely move the total at all---which is the quantitative form of this chapter's qualitative claim that the prize lives in prediction and search, not in automated scientific judgment. The headline is therefore restated once more, and more sharply than the endpoint arithmetic allowed: \emph{a task-weighted central estimate of $\approx$11\%, a likely (10--90\%) range of $\approx$8--18\%, and coherent all-pessimistic/all-optimistic scenario corners near 1\% and 70\%}. The earlier ``3--30\%'' band survives as roughly the outer decile-to-percentile region between the likely range and the corners; the corners of Eq.~\eqref{eq:funnel-simple} survive only if the correlations are assumed away.

What may legitimately be concluded from the lower end, stated in the stock-and-flow discipline this chapter enforces (comparing a present value against an annual industry revenue mixes the two, and the temptation to do so is a standing one): at the pessimistic scenario corner of the task-weighted model---which is $\approx$0.9\%, not the 0.45\% single-pool corner of Eq.~\eqref{eq:funnel-simple}, because the weighted sum cannot put every class at its worst value simultaneously without contradicting the correlations---$0.9\% \times \$3.1\mathrm{T} \approx \$28\mathrm{B}$ per year of incremental research-output flow; and at the 10th-percentile Monte Carlo floor of 7.7\%, $\sim$\$240B per year. \emph{Even the lower end of the modeled uplift represents a socially consequential economic flow, and its present value---applying the Jones--Summers ratio to the increment---exceeds the cost of national AI4S infrastructure by a large multiple.} That is the quantitative form of this chapter's claim. What it does \emph{not} survive is the assumption that acceleration equals discovery, which is why the funnel is drawn.

\subsection{The measured analogue: what national-scale HPC actually returned}
\label{sec:hyperion}

The funnel's parameters are author judgments, and a fair reader will ask whether \emph{any} part of this value chain has ever been measured rather than modeled. One body of evidence exists, and it comes from the adjacent field with three decades of institutional practice: the ROI literature on national high-performance computing, which is the closest measured precedent for what AI4S economics will look like, because leadership HPC is the same asset class---publicly funded computational infrastructure applied to science---one technology generation earlier.

The most complete such study available to this author is Hyperion Research's multi-year assessment of RIKEN's national computing infrastructure and its planned successor~\cite{hyperion-roi}. Two disclosures belong directly beside that sentence: the study was \emph{commissioned by RIKEN}, and the author of this book directs the RIKEN center whose infrastructure it assesses. The reader should weigh everything below accordingly. Its grade, corrected downward from a plain analyst estimate, is \textbf{[A/M]}: what the study reports is not a direct causal measurement of the modeled AI4S quantity but \emph{a commissioned, survey-based, counterfactual economic-impact model informed partly by realized scientific projects}.

Its headline findings: Fugaku's COVID-19 response projects were estimated to have returned \$50--75B to the Japanese economy; 65.7\% of surveyed Fugaku researchers reported already-realized useful results on projects targeting major societal challenges; in an innovation-class mapping of more than 650 HPC projects worldwide, RIKEN's portfolio concentrated in the highest impact-and-importance classes; and the projected returns for the AI-capable successor machine---precisely because AI capability widens the addressable share $a_j$ and the acceleration $s_j$ of Eq.~\eqref{eq:funnel-weighted}---were estimated at \$250B (pessimistic), \$400B+ (realistic), and \$600B+ (optimistic) over its life, two to four times Fugaku's, against an investment two orders of magnitude smaller. The same study put state-directed five-year AI commitments at \$100--200B for China against \$18--27B of US federal spending---the private-public asymmetry that Chapter~\ref{ch:trade} prices---and \$6--10B for Japan.

\begin{casestudy}{Methodology: what is observed, what is modeled, in an HPC ROI study}
Impact studies of this genre chain together components of very different evidentiary strength, and using one honestly requires separating them. \emph{Observed}: that specific projects ran, produced outputs, and reached users---and beneficiary self-reports of usefulness (the 65.7\% figure is this class: real reports, subject to self-selection and social-desirability bias). \emph{Modeled}: the counterfactual (what would have happened without the machine); the attribution of outcomes to the computing infrastructure rather than to the researchers, data, and institutions around it; avoided-loss valuation (the COVID \$50--75B is this class---an estimate of economic collapse \emph{not} suffered, the least verifiable quantity in economics); spillover estimation; and all projected future benefits (the \$250--600B scenarios are projections, full stop). \emph{Unquantified risks}: double counting across projects, displacement (benefits that would have accrued elsewhere), deadweight (projects that needed no subsidy), and the absence of published confidence intervals. A commissioned study has, additionally, a selection gradient in what gets surveyed and how questions are framed. None of this makes the study worthless---it makes it [A/M], and it is why triangulation with independently governed evaluations matters: the UK's ARCHER2 evaluation, commissioned under UKRI's standard evaluation framework with explicit counterfactual scenarios and benefit-cost accounting against public investment, is the best current methodological benchmark for the genre~\cite{archer2-eval}, and a three-way Fugaku--ARCHER2--CERN-class comparison of evaluation methodologies is now on this book's living-appendix work list. The claim this book actually rests on survives the demotion: national-scale computational infrastructure has repeatedly been \emph{evaluated}---by methods far from perfect but more disciplined than anything the enterprise-AI ROI literature offers---at returns that exceed costs by large multiples, concentrated in exactly the task classes where Table~\ref{tab:funnel-classes} puts the value.
\end{casestudy}

What the evidence is good for, at the grade it actually carries: it is \emph{partly ex post} (projects ran; users reported results), it is \emph{independent of this book's thesis} (the analysts were sizing a Japanese procurement, not arguing about China), and its \emph{ratio structure} matches the funnel---the reported value clusters in prediction, simulation, and search applied to societal-scale problems with operational deployment paths, not in the judgment-heavy classes where the funnel assigns near-zero. It is supporting evidence for the funnel's shape and order of magnitude, not a measurement of its value.

\begin{casestudy}{Weather forecasting, end to end through the funnel}
Weather is the one domain where every term of Eq.~\eqref{eq:funnel-weighted} can be observed, which is why it recurs throughout this book and deserves assembly in one place. \emph{Addressability}: the forward model is the dominant computational cost of forecasting, and it is exactly what AI replaced---AIFS operational at ECMWF from February 2025, NOAA's AI system from December 2025 at 0.3\% of the compute of its conventional suite~\cite{ecmwf-aifs,noaa-ai}. \emph{Acceleration}: $10^3$--$10^5\times$ per forecast, measured, not asserted~\cite{aurora,neuralgcm}. \emph{Validation}: forecasting possesses the strongest validation infrastructure in all of science---daily scoring against observation, decades of baseline skill curves---so survival is measured within months. \emph{Deployment}: operational integration inside a year, because national agencies already own the delivery channel. And then, in May 2026, the term everyone forgets: ECMWF discontinued real-time running of all four \emph{external} AI models after an upgrade to its analysis system degraded several of them, while its own AIFS---coupled to its own data assimilation---remained operational~\cite{ecmwf-farewell}. The episode demonstrates five things at once: an AI forecast model is a \emph{component of a production system}, not a free-standing oracle; its performance is contingent on upstream data-assimilation infrastructure it does not control; retraining against infrastructure change is a recurring operational cost, not a one-time expense; deployment requires institutional continuity of exactly the kind Chapter~\ref{ch:consortium} argues must be built rather than rented; and openness was what made the failure \emph{visible}---the retired models could be evaluated, compared, and learned from precisely because they were published. Weather is thus not merely the best AI4S success story; it is the complete systems lesson---architecture, operations, lifecycle cost, trust, and value conversion in one worked example---and the standard against which every other domain's claims in this chapter should be read.
\end{casestudy}

\section{Measuring the scientific frontier: a production metric, not another benchmark}
\label{sec:sci-metric}

Chapter~\ref{ch:model} left the scientific frontier as the one frontier without an accepted metric, and noting a vacancy is worth less than proposing something to fill it. Benchmarks will not serve: a score on a question-answering set measures a component, and the failure ledger above is a catalogue of components that scored well while the system failed. What a ministry actually needs is a \emph{production} metric---value out per resource in:
\begin{equation}
\boldsymbol{\mathcal{S}} \;=\; \frac{\mathbb{E}\left[\mathbf{y}\right]}
{C_{\mathrm{compute}} + C_{\mathrm{human}} + C_{\mathrm{experimental}} + C_{\mathrm{data}} + C_{\mathrm{assurance}}},
\label{eq:sci-metric}
\end{equation}
where the denominator is a scalar cost in currency---$C_{\mathrm{compute}}$ the machine time, $C_{\mathrm{human}}$ the expert labour, $C_{\mathrm{experimental}}$ the wet-lab and instrument cost, $C_{\mathrm{data}}$ acquisition and curation, $C_{\mathrm{assurance}}$ the verification and certification overhead of Chapter~\ref{ch:trust}---and the numerator $\mathbb{E}[\mathbf{y}]$ is the expectation, over the portfolio of projects undertaken, of a \emph{vector} $\mathbf{y}$ of outcome measures. Dividing a vector by a scalar leaves a vector, which is the point: $\boldsymbol{\mathcal{S}}$ is deliberately not collapsed into a single number, because the collapse is where every previous attempt at this metric has smuggled in its assumptions. The components of $\mathbf{y}$: time to independently validated result; expert-hours consumed; replication survival rate; novelty and importance (no single automated novelty metric survives axiomatic scrutiny, and combined metrics do better~\cite{liu2026novelty}); downstream adoption; breadth across domains; false-discovery burden; safety and dual-use externality; and time to operational or regulatory deployment. The denominator's five terms are the cost stack of Chapter~\ref{ch:theory} specialized to science---and the $C_{\mathrm{assurance}}$ term is where the trust stack of Chapter~\ref{ch:trust} enters the economics of discovery directly.

The recent measurement literature supports both the vector form and the humility. The AlphaFold evidence shows \emph{complementarity}---AI-using researchers produced 40\% more novel experimental structures, not fewer---so a metric that scored substitution of human effort would have missed the largest effect~\cite{hill2026,alphafold-impact}. Frontier agents remain below 10\% on demanding scientific literature-discovery tasks, so any metric weighted toward autonomous end-to-end discovery would currently measure noise~\cite{xiong2026}. Scientific value arises from a coupled human--model--simulation--experiment system, and Eq.~\eqref{eq:sci-metric} is written over that system, not over the model. The strategic consequence: whoever operationalizes $\boldsymbol{\mathcal{S}}$ first---runs it across a national research portfolio, publishes the vector, iterates procurement against it---owns the scientific-frontier scoreboard the way MLPerf owns the training scoreboard, and Appendix~\ref{app:dashboard} adds its adoption to the forecast register.

\section{The architecture of the prize}
\label{sec:prize-arch}

The failure ledger and the architecture argument of Part~II point at the same conclusion. AI4S succeeds where a capable model is coupled to simulation, instruments, and data at scale---weather models coupled to assimilation systems, structure prediction coupled to experimental validation, design models coupled to autonomous laboratories. It fails where a model is asked to supply judgement alone. The workload that follows is therefore not chat: it is \emph{long-context agentic orchestration over large batch scientific execution}, which is precisely the two-layer demand structure that Chapter~\ref{ch:sizing} measures, and precisely the workload whose cost is dominated by million-token-context master orchestration rather than by raw dense FLOPs.

Three consequences follow, each connecting to an earlier chapter. \emph{First}, the hardware that serves this best is capacity-rich and bandwidth-balanced rather than FLOP-maximal---the design point of Section~\ref{sec:lpddr6}, and the reason the memory arbitrage is a scientific-infrastructure argument and not only an economic one. \emph{Second}, the models that serve it best are open, because scientific work requires reproducibility, provenance, auditability, and the freedom to fine-tune on data that cannot leave the institution---and because, as Chapter~\ref{ch:consortium} showed, renting the orchestration layer converts a one-seventh workload share into a nine-tenths budget share. \emph{Third}, the failure modes are exactly the ones Chapters~\ref{ch:trust} and~\ref{ch:order66} catalogue: a scientific result produced by an unverifiable artifact through an unauditable chain is not a scientific result, which makes the certification layer a precondition of the prize rather than an accessory to it.

\section{Who is positioned}

The uncomfortable observation, stated at the evidence grade it deserves. China now leads the Nature Index, with output up 22.4\% year on year---the only top-ten country with double-digit growth---and holds the top institution globally along with eight of the top ten in the preceding year's table [V]~\cite{nature-index-2026,nature-index-2025}. Its AI-for-science programme is a decade old at the ministry level and entered the 15th Five-Year Plan alongside fusion~\cite{china-ai4s}. The institutional assets are concrete rather than aspirational: a national scientific agent platform launched by twelve academy institutes with access to 170 million papers and 300-plus computing tools, upgraded a year later on a heterogeneous supercomputing backend~\cite{scienceone,scienceone2}; an autonomous research system wired directly into a national supercomputing network linking thirty-plus centres [P/A]~\cite{scnet}; operational AI weather models at the meteorological administration, one of them open-sourced~\cite{cma-ai}; a commercial AI4S platform serving a thousand universities and 150 corporate clients~\cite{dp-tech}; and---the datapoint that should most concern Western science ministries---a trillion-parameter \emph{open-weight} scientific multimodal foundation model that outperforms the leading closed models by wide margins on scientific reasoning benchmarks~\cite{intern-s1}.

Set against this, the Western response is real and large: a US national mission with an executive-order mandate to double the productivity of American science within a decade, over \$5 billion committed across fifteen-plus agencies and 278 projects, and the two largest scientific AI systems ever procured~\cite{genesis,genesis-scale,genesis-projects}; nineteen European AI factories with a €10 billion envelope and a €20 billion gigafactory programme behind it~\cite{eu-factories,eu-gigafactories}; national systems in the UK and Japan explicitly mandated for AI-driven science~\cite{isambard,fugakunext}. The contest is genuinely joined. But note which side of it is being fought with open weights and cheap tokens, and recall from Chapter~\ref{ch:sizing} that the announced infrastructures on \emph{both} sides underprovision their own research populations by roughly an order of magnitude.

\section{The strategic restatement}

If AI4S is the largest prize, then three of this book's arguments change register.

The token-price argument becomes a \emph{compounding-rate} argument. Chapter~\ref{ch:consortium} showed that purchased frontier tokens at a 20--60$\times$ premium consume three-quarters to nine-tenths of an AI-for-science budget. Read against a 5\%-per-year decline in research productivity and a social return above \$4 per dollar, that is not a procurement inefficiency---it is a tax on the one activity whose returns compound into everything else. A nation that pays it runs its science at a fraction of the rate it could.

The commoditization thesis acquires a beneficiary. Everything Part~II describes---open weights, cheap inference, capacity-first hardware, sovereign-scale clusters---is, from the perspective of this chapter, \emph{good}. Abundant cheap capable AI applied to research is the best available answer to the ideas-getting-harder problem, whoever supplies it. The book's warning is not that commoditization is bad; it is that the country organizing the commons acquires the standards, the adoption, the certification authority, and the scientific compounding that come with it.

And the strategy for everyone else sharpens to a single sentence. If the largest value in AI is scientific, and scientific AI requires abundant compute, open models, verifiable artifacts, and institutions that can carry multi-year continuity, then the correct national investment is not a sovereign chatbot. It is measured capacity (Chapter~\ref{ch:sizing}), pooled internationally (Chapter~\ref{ch:consortium}), running audited open models (Chapter~\ref{ch:trust}) on architectures chosen for long-context agentic science (Chapter~\ref{ch:compute}). That is the whole recommendation of this book, and this chapter is the reason it is worth the money.
